\clearpage
\crshead{ПРИЛОЖЕНИЕ 1 --- Статистические тесты пакета NIST STS}

{\def\LTcaptype{none}%
\begin{longtable}[]{@{}|%
  >{\raggedright\arraybackslash}p{(\linewidth - 6\tabcolsep - 4\arrayrulewidth) * \real{0.0493}}|%
  >{\raggedright\arraybackslash}p{(\linewidth - 6\tabcolsep - 4\arrayrulewidth) * \real{0.3377}}|%
  >{\raggedright\arraybackslash}p{(\linewidth - 6\tabcolsep - 4\arrayrulewidth) * \real{0.6027}}|@{}}
\hline
\begin{minipage}[b]{\hsize}\raggedright
№
\end{minipage} & \begin{minipage}[b]{\hsize}\raggedright
Статистический тест
\end{minipage} & \begin{minipage}[b]{\hsize}\raggedright
Описание теста
\end{minipage} \\ \hline
\endhead
\endlastfoot
1 & Частотный побитовый тест (Frequency Test) & Проверяет баланс между
количеством нулей и единиц в последовательности. Ожидается, что в
достаточно длинной случайной последовательности их соотношение будет
близко к 1:1. \\ \hline
2 & Частотный блочный тест (Frequency Test within a Block) & Частотный
побитовый тест, проводимый в отдельных блоках фиксированного размера \\ \hline
3 & Тест на последовательность одинаковых бит (Runs Test) & Анализирует
чередование нулей и единиц. Проверяет, нет ли слишком длинных серий
подряд идущих одинаковых битов. \\ \hline
4 & Тест на совпадение неперекрывающихся шаблонов (Non-overlapping
Template Matching Test) & Ищет заданные шаблоны в последовательности.
Цель --- выявить ГСП (ГПСП), формирующие слишком часто заданные
непериодические шаблоны \\ \hline
5 & Тест на периодичность (Serial Test) & Данный тест заключается в
подсчете частоты всех возможных перекрываний шаблонов длины~\emph{m}~бит
на протяжении исходной последовательности битов. Целью является
определение действительно ли количество появлений
2\emph{\textsuperscript{m}}~перекрывающихся шаблонов
длиной~\emph{m}~бит, приблизительно такое же как в случае абсолютно
случайной входной последовательности бит. Последняя, как известно,
обладает однообразностью, то есть каждый шаблон длиной~\emph{m}~бит
появляется в последовательности с одинаковой вероятностью. Если
вычисленное в ходе теста значение вероятности~\emph{p}~\textless{} 0,01,
то данная двоичная последовательность не является абсолютно
случайной. \\ \hline
6 & Тест рангов бинарных матриц (Binary Matrix Rank Test) & Разбивает
последовательность на матрицы и вычисляет их ранг. Случайные данные
должны давать ранг, близкий к теоретически ожидаемому. \\ \hline
7 & Тест приблизительной энтропии (Approximate Entropy Test) & Как и
в~тесте на периодичность~в данном тесте акцент делается на подсчёте
частоты всех возможных перекрываний шаблонов длины~\emph{m}~бит на
протяжении исходной последовательности битов. Цель теста --- сравнить
частоты перекрывания двух последовательных блоков исходной
последовательности с длинами~\emph{m}~и~\emph{m+1}~с частотами
перекрывания аналогичных блоков в абсолютно случайной
последовательности. Вычисляемое в ходе теста значение
вероятности~\emph{p}~должно быть не меньше 0,01. В противном случае
(\emph{p}~\textless{} 0,01), двоичная последовательность не является
абсолютно случайной. \\ \hline
8 & Тест кумулятивных сумм (Cumulative Sums Test) & Тест заключается в
максимальном отклонении (от нуля) при произвольном обходе, определяемым
кумулятивной суммой заданных (-1, +1) цифр в последовательности. Цель
данного теста --- определить является ли кумулятивная сумма частичных
последовательностей, возникающих во входной последовательности, слишком
большой или слишком маленькой по сравнению с ожидаемым поведением такой
суммы для абсолютно случайной входной последовательности. \\ \hline
9 & Тест на самую длинную последовательность «1» в блоке (Longest Run of
Ones Test) & Определяет максимальную длину серии единиц в блоках
фиксированного размера. \\ \hline
10 & Тест на совпадение перекрывающихся шаблонов (Overlapping Template
Matching Test) & Аналогичен тесту 4, но шаблоны могут перекрываться.
Полезен для выявления скрытых периодичностей. \\ \hline
11 & Тест на линейную сложность (Linear Complexity Test) & Оценивает,
насколько сложно воспроизвести последовательность с помощью линейного
регистра сдвига. Высокая сложность --- признак хорошей случайности. \\ \hline
12 & Спектральный тест (Discrete Fourier Transform Test) & Преобразует
последовательность в частотную область и ищет пики, которые могут
указывать на периодичность. \\ \hline
13 & Универсальный статистический тест Маурера (Maurer Universal Test) &
Измеряет "сжимаемость" последовательности. Случайные данные не должны
быть алгоритмически сжимаемы. \\ \hline
14 & Тест на произвольные проходы (Random Excursions Test) & Анализирует
случайные блуждания, построенные на основе последовательности.
Проверяет, соответствуют ли они теоретическому распределению. \\ \hline
15 & Тест на произвольные проходы (вариант) (Random Excursions Variant
Test) & Дополняет тест 14, исследуя конкретные состояния блужданий. \\ \hline
\end{longtable}
}
